\documentclass[a4paper,11pt]{article}
\usepackage[utf8]{inputenc}
\usepackage[T1]{fontenc}
\usepackage[french]{babel}
\usepackage[left=1cm,right=1cm,top=.5cm,bottom=0cm]{geometry}
\usepackage{enumitem}
\usepackage{hyperref}
\usepackage{xcolor}
\usepackage{titlesec}
\usepackage{fontawesome5}
\usepackage{lmodern}
\usepackage[normalem]{ulem}

% --- Couleurs RATP ---
\definecolor{primary}{RGB}{0, 100, 60}
\definecolor{secondary}{RGB}{80, 80, 80}

% --- Config Liens ---
\hypersetup{colorlinks=true, linkcolor=primary, filecolor=primary, urlcolor=primary}

% --- Titres ---
\titleformat{\section}{\large\bfseries\color{primary}\uppercase}{}{0em}{}[\titlerule]
\titlespacing{\section}{0pt}{8pt}{5pt}

\begin{document}

% --- EN-TÊTE ---
\begin{center}
    {\Large \textbf{Loïc Aron MBASSI EWOLO}} \\ \vspace{4pt}
    \textbf{\large Alternant ML Engineer IA Générative} \\
    \textit{\small Disponible dès septembre 2026} \\ \vspace{4pt}
    \small
    \faMapMarker\ Cergy, France (Mobile IDF) \quad
    \faPhone\ (+33) 7 59 52 33 98 \quad
    \faEnvelope\ \href{mailto:loicaron.mbassiewolo@ensea.fr}{loicaron.mbassiewolo@ensea.fr} \\ \vspace{2pt}
    \faLinkedin\ \href{https://www.linkedin.com/in/loic-mbassi/}{linkedin.com/in/loic-mbassi} \quad
    \faGithub\ \href{https://github.com/Nameless0l}{github.com/Nameless0l} \quad
    \faGlobe\ \href{https://mbassiloic.tech}{mbassiloic.tech}
\end{center}

\vspace{-6pt}

% --- PROFIL ---
\noindent\small{Orchestration de 5 agents IA via LangChain et Google ADK avec pipeline RAG, déploiement FastAPI/Docker, et NLP à grande échelle (BERT/Transformers) à l'INRIA Saclay : ces expériences me préparent à contribuer directement au développement et à l'industrialisation de vos solutions d'IA Générative au sein de la Data Factory du groupe RATP.}

% --- FORMATION ---
\section{Formation}
\begin{itemize}[leftmargin=*, label=\tiny$\bullet$, nosep]
    \item \textbf{ENSEA -- École Nationale Supérieure de l'Électronique et de ses Applications} \hfill \textit{2025 -- 2027} \\
    \small{Double diplôme d'Ingénieur -- Systèmes Embarqués, Signal, IA. Cursus : Deep Learning, Traitement du Signal.}
    \item \textbf{École Polytechnique de Yaoundé (1\textsuperscript{ère} école d'ingénieur CEMAC)} \hfill \textit{2021 -- 2025} \\
    \small{Diplôme d'Ingénieur de Conception (Master 2) : \textbf{Mathématiques Appliquées}, Génie Logiciel, Algorithmique.}
\end{itemize}

% --- COMPÉTENCES ---
\section{Compétences Techniques}
\begin{itemize}[leftmargin=*, nosep]
    \item \small{\textbf{IA Générative \& Agents :} LangChain, Google ADK, RAG (vectorisation, chunking, prompting), Gemini, GPT-2, BERT/Transformers, orchestration multi-agents, prompt engineering.}
    \item \small{\textbf{ML \& NLP :} Python, PyTorch, TensorFlow, scikit-learn, pandas, NumPy, traitement du langage naturel, séries temporelles.}
    \item \small{\textbf{Développement \& Cloud :} FastAPI, APIs REST, Docker, Git, Azure (notions), Poetry, CI/CD. Qualité de code (linting, pre-commit).}
    \item \small{\textbf{Langues :} Français (natif), Anglais (professionnel / technique).}
\end{itemize}

% --- EXPÉRIENCES / PROJETS ---
\section{Expériences \& Projets}
\begin{itemize}[leftmargin=*, label={-}]

\item \textbf{Système Multi-Agents Agricole -- Orchestration LLM \& RAG} \hfill \textit{2024 -- 2025} \\
\textit{Projet de recherche} \hfill \textit{Python, LangChain, Google ADK, Gemini, RAG, FastAPI, Docker}
\vspace{-2pt}
    \begin{itemize}[leftmargin=0.4cm, nosep, label=\tiny$\bullet$]
        \item \small{Conception et déploiement de 5 agents IA collaboratifs (Google ADK) utilisant Gemini comme moteur LLM, pour des recommandations agricoles multi-critères (météo, cultures, marchés).}
        \item \small{Pipeline RAG complet : vectorisation des connaissances, recherche sémantique par similarité cosinus, augmentation contextuelle du prompt, évaluation de la pertinence des réponses.}
        \item \small{Orchestration d'APIs externes (OpenWeather, marchés), résolution de conflits entre agents, déploiement production via FastAPI/Docker (Poetry).}
    \end{itemize}
\vspace{3pt}

\item \textbf{Stage Recherche -- INRIA Saclay, Équipe TRIBE (PEPR MOBIDEC)} \hfill \textit{Avril -- Août 2026} \\
\textit{Institut Polytechnique de Paris} \hfill \textit{Python, NLP, BERT/Transformers, pandas, scikit-learn}
\vspace{-2pt}
    \begin{itemize}[leftmargin=0.4cm, nosep, label=\tiny$\bullet$]
        \item \small{Pipeline d'analyse NLP à grande échelle des politiques de confidentialité d'applications mobiles (BERT/Transformers) : évaluation transparence, conformité RGPD.}
        \item \small{Conception d'un score de confiance interprétable synthétisant les résultats d'analyse. Contribution à des publications scientifiques.}
    \end{itemize}
\vspace{3pt}

\item \textbf{Assistant de Recherche -- MILA (Institut Québécois d'IA)} \hfill \textit{Juin -- Sept 2025} \\
\textit{Remote} \hfill \textit{Python, PyTorch, interprétabilité des modèles}
\vspace{-2pt}
    \begin{itemize}[leftmargin=0.4cm, nosep, label=\tiny$\bullet$]
        \item \small{Étude du phénomène de Grokking (généralisation tardive des réseaux de neurones) sur MLP : reproduction d'expériences état de l'art, analyse convergence mathématique.}
        \item \small{Visualisation des dynamiques d'entraînement, étude du comportement émergent des réseaux profonds.}
    \end{itemize}
\vspace{3pt}

\item \textbf{Snappy -- Plateforme de Messagerie Intelligente (Chatbots RAG)} \hfill \textit{2024 -- 2025} \\
\textit{Projet académique 4\textsuperscript{e} année ENSPY} \hfill \textit{Spring Boot, Python/Flask, RAG, Gemini, Redis, Docker}
\vspace{-2pt}
    \begin{itemize}[leftmargin=0.4cm, nosep, label=\tiny$\bullet$]
        \item \small{J'ai participé à la conception de chatbots intelligents avec RAG (Gemini Flash/Pro + vector store) pour réponses contextuelles à partir de documents d'organisation.}
        \item \small{Recherche sémantique par embeddings vectoriels (similarité cosinus), enrichissement des réponses avec le contexte documentaire pertinent.}
    \end{itemize}
\vspace{3pt}

\end{itemize}

% --- DISTINCTIONS ---
\section{Distinctions \& Leadership}
\begin{itemize}[leftmargin=*, nosep, label=\tiny$\bullet$]
    \item \small{\textbf{1\textsuperscript{er} Prix} -- IndabaX Africa 2025 (IA Santé) : nettoyage de données médicales, déploiement API Flask/Streamlit.}
    \item \small{\textbf{1\textsuperscript{er} Prix} -- CITS National 2024 : pipeline ML optimisation de routes, réduction 20\% des coûts, 75 équipes.}
    \item \small{\textbf{Lead AI Cell ENSPY} : animation workshops ML, fine-tuning LLM. Collaboration Google DeepMind. Articles techniques sur Medium.}
\end{itemize}

\end{document}
